% The shared design of the three PDFs. One file, because three drifting copies of a
% typographic system is the same failure as three copies of a constant.
%
% lualatex, for fontspec. TeX Gyre Pagella (Palatino) for text, Heros (Helvetica) for
% headings and figure labels, Cursor for code. The charts are matplotlib PDFs set in the
% same two faces, so a figure and the paragraph under it are the same typeface.
% pdflatex, not lualatex: texlive-luatex is not installed here, so fontspec has no font
% loader — and a document that only builds after someone installs a package is a document
% that does not build. The TeX Gyre Type1 faces ship with texlive-fonts-recommended and give
% the same typefaces: Pagella (Palatino) for text, Heros (Helvetica) for headings and figure
% labels, Cursor for code. The matplotlib charts are set in the same two, so a figure and the
% paragraph under it share a typeface.
\usepackage[T1]{fontenc}
\usepackage[utf8]{inputenc}
\usepackage{tgpagella}
\usepackage{tgheros}
\usepackage{tgcursor}
\renewcommand{\sfdefault}{qhv}
\renewcommand{\ttdefault}{qcr}
\usepackage{textcomp}
\usepackage{amssymb}
\usepackage{geometry}
\usepackage{microtype}
\usepackage{booktabs}
\usepackage{longtable}
\usepackage{array}
\usepackage{tabularx}
\usepackage[table]{xcolor}
\usepackage{titlesec}
\usepackage{enumitem}
\usepackage{graphicx}
\usepackage{caption}
\usepackage{fancyhdr}
\usepackage{tikz}
\usepackage[hidelinks]{hyperref}

% Old-style figures are Pagella's default in text; \lining forces lining figures where a
% column of numbers has to align.
\newcommand{\lining}[1]{{\fontfamily{qpl}\selectfont #1}}

% Characters the sources use that T1 does not carry.
\DeclareUnicodeCharacter{2212}{--}
\DeclareUnicodeCharacter{2264}{$\leq$}
\DeclareUnicodeCharacter{2265}{$\geq$}
\DeclareUnicodeCharacter{2192}{$\rightarrow$}
\DeclareUnicodeCharacter{00B7}{\textperiodcentered}
\DeclareUnicodeCharacter{00A2}{\textcent}

\definecolor{ink}{HTML}{16161A}
\definecolor{muted}{HTML}{55555F}
\definecolor{faint}{HTML}{8A8A94}
\definecolor{rule}{HTML}{D8D8DE}
\definecolor{accent}{HTML}{C4185C}      % the estate pink, darkened for paper
\definecolor{accentpale}{HTML}{FBEEF3}
\definecolor{cool}{HTML}{2C5B87}
\definecolor{coolpale}{HTML}{EEF3F8}
\definecolor{warnpale}{HTML}{FFF6E8}
\definecolor{warnrule}{HTML}{D9A441}

\color{ink}
\hypersetup{colorlinks=true, linkcolor=cool, urlcolor=cool, citecolor=cool}

\titleformat{\section}{\sffamily\bfseries\large\color{ink}}{\thesection}{0.7em}{}
\titleformat{\subsection}{\sffamily\bfseries\normalsize\color{ink}}{\thesubsection}{0.6em}{}
\titleformat{\subsubsection}{\sffamily\bfseries\small\color{muted}}{}{0em}{}
\titlespacing*{\section}{0pt}{1.7\baselineskip}{0.6\baselineskip}
\titlespacing*{\subsection}{0pt}{1.2\baselineskip}{0.4\baselineskip}
\titlespacing*{\subsubsection}{0pt}{1.0\baselineskip}{0.25\baselineskip}

\setlist[itemize]{leftmargin=1.2em, itemsep=0.25em, topsep=0.4em, parsep=0pt}
\setlist[enumerate]{leftmargin=1.4em, itemsep=0.25em, topsep=0.4em, parsep=0pt}

\captionsetup{font={sf,small,color=muted}, labelfont={sf,bf,color=ink},
              labelsep=period, justification=raggedright, singlelinecheck=false,
              skip=6pt}

\setlength{\parindent}{0pt}
\setlength{\parskip}{0.46\baselineskip}
\renewcommand{\arraystretch}{1.16}
\setlength{\tabcolsep}{7pt}
\arrayrulecolor{rule}   % needs [table]{xcolor}, above

% A rule the width of the text, used to close a section quietly.
\newcommand{\thinrule}{\par\vspace{0.3\baselineskip}%
  {\color{rule}\hrule height 0.4pt}\par\vspace{0.3\baselineskip}}

% ---- callouts -------------------------------------------------------------------------
% Three registers, and the colour says which: accent for the claim a page is making, cool
% for a fact worth pausing on, amber for a limit the reader must carry forward. A slab is a
% tinted minipage under a coloured rule — no tcolorbox on this machine, and a plain
% construction is one less thing that can break a build on someone else's.
\newcommand{\slab}[4]{% {bg}{rulecolour}{TITLE}{body}
  \par\medskip\noindent
  {\color{#2}\hrule height 1.2pt}%
  \colorbox{#1}{\begin{minipage}{\dimexpr\linewidth-2\fboxsep\relax}
    \vspace{2pt}
    {\sffamily\bfseries\footnotesize\color{#2}\MakeUppercase{#3}}\par
    \vspace{0.15\baselineskip}
    #4
    \vspace{2pt}
  \end{minipage}}%
  \par\medskip}
\newcommand{\headline}[2]{\slab{accentpale}{accent}{#1}{#2}}
\newcommand{\keypoint}[2]{\slab{coolpale}{cool}{#1}{#2}}
\newcommand{\limitbox}[2]{\slab{warnpale}{warnrule}{#1}{#2}}

% ---- the figure store -----------------------------------------------------------------
% \F{P10000.cycle.tph} resolves a key from research/figures.json, which is regenerated from
% the live model. A key that does not exist is a build error, not a blank.
\makeatletter
\newcommand{\F}[1]{%
  \ifcsname fig@#1\endcsname\csname fig@#1\endcsname
  \else\GenericError{}{Figure key '#1' is not in figures.tex}{}{}\fi}
\makeatother

% ---- figures, placed rather than floated ------------------------------------------------
% LaTeX floats are for papers where the figure may go anywhere. In a designed report the
% figure belongs beside the sentence that introduces it, and letting it drift left half-empty
% pages and pushed a five-page brief off the end. So a figure is an ordinary block: it appears
% exactly where it was written, numbered by hand.
\newcounter{plainfig}
\newcommand{\placedfig}[3]{% {path}{caption}{width fraction}
  \par\addvspace{0.9\baselineskip}\noindent
  \begin{center}\includegraphics[width=#3\linewidth]{#1}\end{center}
  \vspace{-0.2\baselineskip}
  \refstepcounter{plainfig}%
  {\sffamily\small\color{muted}\textbf{\color{ink}Figure \theplainfig.} #2\par}
  \addvspace{0.9\baselineskip}}

% Big number, for the statistic slabs.
\newcommand{\bignum}[2]{{\sffamily\bfseries\Large\color{accent}#1}\,{\sffamily\small\color{muted}#2}}
